% Môn: Vật lý
% Lớp: 12
% Chương: Chương II. Khí lí tưởng
% Bài: L12C2 Bài 9. Định luật Boyle · Bài toán về quá trình đẳng nhiệt
% ID: L12C1B2-01-DS
% Mức: NB
% ID: L12C1B2-01-DS


% ID: L12C1B2-01-TLN
% Mức: VD
% ID: L12C1B2-01-TLN
% ID: L12C1B2-01-TLN
% Mức: VD
% ID: L12C2B9-57-TLN
%=== Câu 3 ===%
\dangbt{Bài toán về quá trình đẳng nhiệt}
\begin{ex}
% ID: L12C2B9-57-TLN
% Mức: VD
Một quả bóng có thể tích khi căng là $5,6\,\ell$. Mỗi lần bơm đưa vào bóng $0,32\,\ell$ không khí ở áp suất $1\,\text{atm}$. Ban đầu bóng không có không khí. Coi nhiệt độ không đổi. Sau $35$ lần bơm, áp suất khí trong bóng là bao nhiêu atm?
\shortans{$2,0$}
\loigiai{
Áp dụng định luật Boyle cho toàn bộ lượng khí được bơm vào:
$p_0 n\Delta V=pV.Suy ra:
 p=\dfrac{p_0 n\Delta V}{V}$ = $\dfrac{1\cdot 35\cdot 0,32}{5,6}=2,0\text{atm}$. 
Vậy áp suất khí trong bóng là $2,0\,\text{atm}$.
}
\end{ex}
